\chapter*{ТЕРМИНЫ И ОПРЕДЕЛЕНИЯ}
\addcontentsline{toc}{chapter}{ТЕРМИНЫ И ОПРЕДЕЛЕНИЯ}
\begingroup
\setlength{\parindent}{0pt}
% Перечень располагается столбцом, слева без абзацного отступа в алфавитном порядке,
% справа через тире — определение, без знаков препинания в конце строки (п. 4.9.2)
% Термины пишутся без жирного шрифта
Граничная верификация модели --- проверка имитационной модели на удовлетворение обязательным свойствам в граничных значениях её параметров и в режимах редукции\\[0.5em]
Косинусная близость --- мера сходства двух векторов, определяемая как косинус угла между ними\\[0.5em]
Латинский гиперкуб --- план эксперимента, обеспечивающий равномерное стратифицированное покрытие пространства параметров заданным числом точек при сокращении объёма по сравнению с полным факторным перебором\\[0.5em]
Общие случайные числа --- приём планирования имитационных экспериментов, при котором сравниваемые альтернативы выполняются на одной и той же реализации источников случайности для снижения шума попарного сравнения\\[0.5em]
Политика рекомендаций --- модуль, отображающий программу конференции и профиль участника в упорядоченный список рекомендуемых докладов в текущем временном слоте\\[0.5em]
Социальное заражение --- компонент модели поведения участника, отражающий монотонное усиление склонности к выбору доклада с ростом числа уже сделанных в когорте выборов того же доклада\\[0.5em]
Сценарный анализ --- подход к принятию решений в условиях глубокой неопределённости, в котором стратегии оцениваются на множестве правдоподобных сценариев, а не на единственной точечной прогнозной модели\\[0.5em]
Функция полезности --- числовая функция, по значениям которой формируется распределение вероятностей выбора участника между доступными в слоте докладами
\endgroup
